\subsection{Extensivity}
A category is extensive if 
it has all finite coproducts and for all $X_1,X_2\in\catc$ the functor
\[
   F_+ : \catc / X_1 \times \catc / X_2 \to \catc /(X_1+X_2)
\]
is an equivalence of categories. We define $F_+$ as follows
\[
    F_+((A_1,f_1: A_1 \to X_1),(A_2,f_2: A_2 \to X_2)) \mapsto 
    (A_1\times A_2, (\iota_{X_1}\circ f_1) + (\iota_{X_2}\circ f_2))
\]
where $(\iota_{X_1}\circ f_1) + (\iota_{X_2}\circ f_2)$ is the unique map 
(existence and uniqueness follows from universal property of $A_1+A_2$) such that the following diagram commutes 
\[\begin{tikzcd}
	{A_1} & {A_1+A_2} & {A_2} \\
	{X_1} & {X_1+X_2} & {X_2}
	\arrow["{\iota_{A_1}}", from=1-1, to=1-2]
	\arrow["{f_1}"', from=1-1, to=2-1]
	\arrow["{\exists !}", dashed, from=1-2, to=2-2]
	\arrow["{\iota_{A_2}}"', from=1-3, to=1-2]
	\arrow["{f_2}", from=1-3, to=2-3]
	\arrow["{\iota_{X_1}}"', from=2-1, to=2-2]
	\arrow["{\iota_{X_2}}", from=2-3, to=2-2]
\end{tikzcd}\]
Given a morphism in $\catc / X_1 \times \catc / X_2$
\[
    (h_1,h_2) : ((A_1,f_1: A_1 \to X_1),(A_2,f_2: A_2 \to X_2)) \to 
    ((B_1,g_1: B_1 \to X_1),(B_2,g_2: B_2 \to X_2))
\]
so $h_1 : A_1 \to B_1 \land h_2 : A_2 \to B_2$ such that the following 
two triangles commute
\[\begin{tikzcd}
	{A_1} & {B_1} & {A_2} & {B_2} \\
	& {X_1} && {C_2}
	\arrow["{h_1}", from=1-1, to=1-2]
	\arrow["{f_1}"', from=1-1, to=2-2]
	\arrow["{g_1}", from=1-2, to=2-2]
	\arrow["{h_2}", from=1-3, to=1-4]
	\arrow["{f_2}"', from=1-3, to=2-4]
	\arrow["{g_2}", from=1-4, to=2-4]
\end{tikzcd}\]
we define $F_+(h_1,h_2)$ as the unique morphism $(\iota_{B_1}\circ h_1)+(\iota_{B_2}\circ h_2)$
which make the following diagram commute
\[\begin{tikzcd}
	{A_1} & {A_1+A_2} & {A_2} \\
	{B_1} & {B_1+B_2} & {B_2}
	\arrow["{\ins{A_1}}", from=1-1, to=1-2]
	\arrow["{h_1}", from=1-1, to=2-1]
	\arrow["{\exists !}", dashed, from=1-2, to=2-2]
	\arrow["{\ins{A_2}}"', from=1-3, to=1-2]
	\arrow["{h_2}", from=1-3, to=2-3]
	\arrow["{\ins{B_1}}", from=2-1, to=2-2]
	\arrow["{\ins{B_2}}"', from=2-3, to=2-2]
\end{tikzcd}\]
Now to show that it is a morphism in the slice category,
we have to verify that the following triangle commutes
\[\begin{tikzcd}
	{A_1+A_2} && {B_1+B_2} \\
	\\
	&& {X_1+X_2}
	\arrow["{(\ins{B_1}\circ h_1)+(\ins{B_2}\circ h_2)}", from=1-1, to=1-3]
	\arrow["{(\ins{X_1}\circ f_1)+(\ins{X_2}\circ f_2)}"', from=1-1, to=3-3]
	\arrow["{(\ins{X_1}\circ g_1)+(\ins{X_2}\circ g_1)}", from=1-3, to=3-3]
\end{tikzcd}\]
\begin{lemma}
    Morphism from coproduct is uniquely determined by
    composition with its insertions. Consider the following diagram
\[\begin{tikzcd}
	{X_1} & {X_1+X_2} & {X_2} \\
	& X
	\arrow["{\ins{X_1}}", from=1-1, to=1-2]
	\arrow["{g\circ\ins{X_1}=f\circ\ins{X_1}}"', from=1-1, to=2-2]
	\arrow["f", shift left, from=1-2, to=2-2]
	\arrow["g"', shift right, from=1-2, to=2-2]
	\arrow["{\ins{X_2}}"', from=1-3, to=1-2]
	\arrow["{f\circ\ins{X_2}=g\circ\ins{X_2}}", from=1-3, to=2-2]
\end{tikzcd}\]
then $f=g$.
\end{lemma}
\begin{proof}[Proof]
    Follows directly from the universal property of coproducts that
    there exists a unique arrow $\kappa : X_1\times X_2\to X$ that makes
	the triangles commute, but both $f$ and $g$ make the triangles commute,  
    so $f=\kappa=g$.
\end{proof}
Consider the diagram
\[\begin{tikzcd}
	{A_1} && {A_1+A_2} && {A_2} \\
	{B_1} && {B_1+B_2} && {B_2} \\
	{X_1} && {X_1+X_2} && {X_2}
	\arrow["{\ins{A_1}}", color={rgb,255:red,214;green,92;blue,92}, from=1-1, to=1-3]
	\arrow["{f_1}", from=1-1, to=2-1]
	\arrow["{h_1}"{description}, color={rgb,255:red,214;green,92;blue,92}, curve={height=24pt}, from=1-1, to=3-1]
	\arrow[from=1-3, to=2-3]
	\arrow[color={rgb,255:red,214;green,92;blue,92}, curve={height=-24pt}, from=1-3, to=3-3]
	\arrow["{\ins{A_2}}"', color={rgb,255:red,214;green,92;blue,92}, from=1-5, to=1-3]
	\arrow["{f_2}"', from=1-5, to=2-5]
	\arrow["{h_2}"{description}, color={rgb,255:red,214;green,92;blue,92}, curve={height=-24pt}, from=1-5, to=3-5]
	\arrow["{\ins{B_1}}", from=2-1, to=2-3]
	\arrow["{g_1}", from=2-1, to=3-1]
	\arrow[from=2-3, to=3-3]
	\arrow["{\ins{B_2}}"', from=2-5, to=2-3]
	\arrow["{g_2}"', from=2-5, to=3-5]
	\arrow["{\ins{X_1}}", color={rgb,255:red,214;green,92;blue,92}, from=3-1, to=3-3]
	\arrow["{\ins{X_2}}"', color={rgb,255:red,214;green,92;blue,92}, from=3-5, to=3-3]
\end{tikzcd}\]
To prove that the inner triangle (whose unmarked arrows are as above) commutes
we chase around the diagram starting from the insertions by the lemma above.
We know by assumption that the leftmost and rightmost 
triangles commute and that all the four black squares
commute by definition of the unmarked arrows and the
same for the two red squares.\par
Now we ought to show that $F_+$ really is a functor. We have
\[
    F_+(1_{A_1},1_{A_2}) = (\iota_{A_1}\circ 1_{A_1})+(\iota_{A_2}\circ 1_{A_2}) = \iota_{A_1}+\iota_{A_2} = 1_{A_1+A_2} = 1_{F_+((A_1,f_1),(A_2,f_2))}
\]
the penultimate equality again follows by the universal property of $A_1+A_2$.\par
Given two morphisms $(h_1: A_1\to B_1,h_2: A_2\to B_2),(k_1:B_1\to C_1 ,k_2: B_2\to C_2)$ we have
\[
    F_+((k_1,k_2)\circ(h_1,h_2)) = F_+((k_1\circ h_1,k_2\circ h_2))
    = (\iota_{C_1}\circ k_1\circ h_1)+(\iota_{C_2}\circ k_2\circ h_2)
\]
Notice that to prove this equals the desired
\[((\iota_{C_1}\circ k_1)+(\iota_{C_2}\circ k_2))\circ((\ins{B_1}\circ h_1)+(\ins{B_2}\circ h_2))\] 
amounts to showing that the unmarked inner triangle in the following diagram commutes provided
that the four square diagrams commute (the square once again commute because of the universality of coproduct)
\[\begin{tikzcd}
	{A_1} && {A_1+A_2} && {A_2} \\
	{B_1} && {B_1+B_2} && {B_2} \\
	{C_1} && {C_1+C_2} && {C_2}
	\arrow["{\iota_{A_1}}", from=1-1, to=1-3]
	\arrow["{h_1}"', from=1-1, to=2-1]
	\arrow[from=1-3, to=2-3]
	\arrow[curve={height=-24pt}, from=1-3, to=3-3]
	\arrow["{\iota_{A_2}}"', from=1-5, to=1-3]
	\arrow["{k_1}", from=1-5, to=2-5]
	\arrow["{\iota_{B_1}}", from=2-1, to=2-3]
	\arrow["{h_2}"', from=2-1, to=3-1]
	\arrow[from=2-3, to=3-3]
	\arrow["{\iota_{B_2}}"', from=2-5, to=2-3]
	\arrow["{k_2}", from=2-5, to=3-5]
	\arrow["{\iota_{C_1}}", from=3-1, to=3-3]
	\arrow["{\iota_{C_2}}"', from=3-5, to=3-3]
\end{tikzcd}\]
we again make use of the lemma and chase around the diagram until we get the equality. 
It follows, that $F_+((k_1,k_2)\circ(h_1,h_2)) = F_+((k_1,k_2))\circ F_+(h_1,h_2)$.